\section{Thiết kế Kiến trúc và Cơ chế Hệ thống (V2)}

\subsection{Tổng quan Kiến trúc Modular Monolith}
Để đáp ứng đồng thời yêu cầu về khả năng mở rộng chức năng, tính toàn vẹn dữ liệu và khả năng kiểm toán truy vết, Evidence Atlas được thiết kế theo mô hình **Khối nguyên khối mô-đun hóa (Modular Monolith)** \src{1}. 

Thay vì phân tán thành nhiều vi dịch vụ (\textit{microservices}) phức tạp gây độ trễ mạng và khó khăn trong việc duy trì giao dịch phân tán, toàn bộ các ranh giới nghiệp vụ của hệ thống được đóng gói trong các mô-đun Python độc lập, chia sẻ một cơ sở dữ liệu quan hệ bền vững **PostgreSQL** duy nhất.

\begin{figure}[htbp]
\centering
\begin{tikzpicture}[node distance=1.4cm, auto, >=stealth, every node/.style={font=\small}]
    \tikzstyle{client} = [rectangle, draw, fill=gray!10, text width=9cm, text centered, rounded corners, minimum height=0.8cm, draw=gray!60]
    \tikzstyle{gateway} = [rectangle, draw, fill=blue!15, text width=9cm, text centered, rounded corners, minimum height=0.9cm, draw=blue!60, font=\bfseries\small]
    \tikzstyle{supervisor} = [rectangle, draw, fill=green!10, text width=9cm, text centered, rounded corners, minimum height=1.1cm, draw=green!60]
    \tikzstyle{experts} = [rectangle, draw, fill=yellow!15, text width=2.1cm, text centered, rounded corners, minimum height=1cm, draw=yellow!60, font=\footnotesize]
    \tikzstyle{storage} = [cylinder, shape border rotate=90, draw, fill=orange!15, text width=2.5cm, text centered, minimum height=1cm, draw=orange!60, font=\footnotesize]
    \tikzstyle{line} = [draw, -latex', line width=0.8pt]

    \node [client] (ui) {Khách hàng Web / Ứng dụng ngoài (X-API-Key, JSON / SSE)};
    \node [gateway, below of=ui] (gw) {Agent Gateway (Xác thực, Phân quyền Role-based, Giới hạn Budget)};
    \node [supervisor, below of=gw] (sup) {\textbf{V2 Supervisor (Bounded Multi-Expert Orchestration)} \\ Lập kế hoạch, Gọi tuần tự, Exact-Span Matching, Tổng hợp câu trả lời};
    
    \node [experts, below of=sup, xshift=-3.3cm] (exp1) {\textbf{Product} \\ Lọc / Xếp hạng};
    \node [experts, below of=sup, xshift=-1.1cm] (exp2) {\textbf{Review} \\ Trích xuất ý kiến};
    \node [experts, below of=sup, xshift=1.1cm] (exp3) {\textbf{Trust} \\ Phân tích phàn nàn};
    \node [experts, below of=sup, xshift=3.3cm] (exp4) {\textbf{Market} \\ Thống kê danh mục};

    \node [storage, below of=exp2, xshift=-1cm, yshift=-0.5cm] (pg1) {PostgreSQL \\ Hội thoại \& Lượt};
    \node [storage, below of=exp3, xshift=1cm, yshift=-0.5cm] (pg2) {PostgreSQL \\ Snapshot \& Kho RAG};

    \path [line] (ui) -- (gw);
    \path [line] (gw) -- (sup);
    \path [line] (sup) -- (exp1);
    \path [line] (sup) -- (exp2);
    \path [line] (sup) -- (exp3);
    \path [line] (sup) -- (exp4);
    \path [line] (exp1) -- (pg2);
    \path [line] (exp2) -- (pg2);
    \path [line] (exp3) -- (pg2);
    \path [line] (exp4) -- (pg2);
    \path [line] (sup) -- (pg1);
\end{tikzpicture}
\caption{Kiến trúc tổng thể hệ thống Evidence Atlas (Phiên bản API v2).}
\label{fig:v2_architecture}
\end{figure}

Giao diện lập trình ứng dụng HTTP `/api/v2` là điểm giao tiếp duy nhất của hệ thống, hỗ trợ hai phương thức tương tác:
\begin{itemize}
    \item \textbf{POST JSON tiêu chuẩn:} Nhận câu hỏi và trả về toàn bộ kết quả sau khi quy trình tổng hợp hoàn tất.
    \item \textbf{Server-Sent Events (SSE) Streaming:} Truyền phát các sự kiện tiến trình theo thời gian thực (bắt đầu phân tích, gọi chuyên gia, tìm kiếm dữ liệu, trích dẫn bằng chứng) và dòng câu trả lời hoàn thiện, kết thúc bằng duy nhất một sự kiện thiết bị đầu cuối (\textit{terminal event}).
\end{itemize}

\subsection{Cổng Tác tử (Agent Gateway) và An toàn Phân quyền}
Để ngăn chặn các hành vi leo thang đặc quyền (\textit{privilege escalation}) hoặc can thiệp trái phép vào dữ liệu, toàn bộ các yêu cầu từ phía máy khách đều phải đi qua \textbf{Agent Gateway} trước khi tiếp cận tầng điều phối nghiệp vụ \src{1}.

Agent Gateway hoạt động theo nguyên tắc đóng kín khi có nghi vấn (\textit{fail-closed}):
\begin{enumerate}
    \item \textbf{Xác thực định danh qua khóa API (Principal Authentication):} Khóa API gửi qua tiêu đề HTTP \texttt{X-API-Key} được ánh xạ với một định danh chủ thể (\textit{principal}) cụ thể.
    \item \textbf{Phân quyền theo phạm vi và vai trò (Role-Based Access Control):}
    \begin{itemize}
        \item \textbf{Vai trò Người mua (Shopper):} Chỉ được cấp quyền đọc dữ liệu danh mục công khai (\texttt{ecommerce.read}) và thực hiện các hành động tư vấn, tìm kiếm hoặc chuẩn bị giỏ hàng ảo mang tính đề xuất.
        \item \textbf{Vai trò Người bán (Merchant):} Được cấp quyền đọc báo cáo phân tích kho hàng và đề xuất điều chỉnh thông tin bán hàng.
    \end{itemize}
    \item \textbf{Kiểm soát ủy quyền hành động (Action Authorization):} Mọi hành động thao tác dữ liệu (như ghi đơn hàng hay sửa đổi giá) đều được tạo dưới dạng bản ghi đề xuất (\textit{Proposal}) và yêu cầu xác nhận rõ ràng, tuyệt đối không cho phép mô hình ngôn ngữ tự ý thực thi các thao tác ghi dữ liệu mà không có chữ ký của người dùng.
\end{enumerate}

Nhờ cơ chế kiểm soát nghiêm ngặt này, trong toàn bộ 720 lượt thử nghiệm thuộc bộ benchmark P8, hệ thống ghi nhận \textbf{0 vi phạm phân quyền} (tỷ lệ ủy quyền hợp lệ đạt $100{,}0\%$).

\subsection{Bộ điều phối Bounded Hybrid Multi-Expert Supervisor}
Khác với các hệ thống Multi-Agent đàm phán tự do, Evidence Atlas sử dụng bộ điều phối \textbf{Supervisor} với quy trình lập kế hoạch có chặn trên (\textit{bounded execution}) nhằm đảm bảo tính tiền định và kiểm soát độ trễ \src{1}.

\subsubsection{Trách nhiệm của 4 Domain Experts}
Hệ thống cố định 4 chuyên gia chức năng độc lập:
\begin{itemize}
    \item \texttt{product\_agent}: Tiếp nhận tiêu chí lọc (thể loại, khoảng giá, số lượng), truy vấn bảng cơ sở dữ liệu danh mục sách và tính toán điểm xếp hạng ứng viên dựa trên độ tương đồng và mức độ phổ biến.
    \item \texttt{review\_agent}: Trích xuất các bản ghi nhận xét của người mua gắn liền với danh sách sách đã được chọn, tổng hợp các ý kiến khen chê về nội dung văn học và trải nghiệm đọc.
    \item \texttt{trust\_agent}: Phân tích các dấu hiệu bất thường định tính, nhận diện các phàn nàn phổ biến liên quan đến bao bì, giao hàng trễ, hoặc sách móp méo bằng các mẫu quy tắc được phiên bản hóa.
    \item \texttt{market\_agent}: Tính toán các chỉ số thống kê cắt ngang tổng thể trên tập snapshot (như giá trung bình theo tác giả, tỷ lệ phân bổ sao) để cung cấp bối cảnh toàn cảnh.
\end{itemize}

\subsubsection{Quy trình điều phối và Ràng buộc thực thi}
Khi tiếp nhận câu hỏi của người dùng, Supervisor thực thi tuần tự theo chu trình có giới hạn nghiêm ngặt được định nghĩa trong hợp đồng dữ liệu \texttt{app/v2/contracts.py}:

\begin{figure}[htbp]\small
\centering
\begin{tabularx}{\linewidth}{lP{28mm}Y}\toprule
Ràng buộc & Giá trị giới hạn & Ý nghĩa kỹ thuật\\\midrule
\texttt{MAX\_CANDIDATES} & 5 cuốn sách & Số lượng ứng viên tối đa được đưa vào vòng đánh giá chi tiết\\
\texttt{MAX\_PLAN\_REVISIONS} & 1 lần & Số lần tối đa kế hoạch được phép điều chỉnh sau khi xem dữ chứng chặng 1\\
\texttt{MAX\_ADDED\_READ\_STEPS} & 2 bước & Số lượng bước đọc bổ sung tối đa được phép chèn thêm vào vòng lặp\\
\texttt{MAX\_EXPERT\_STEPS} & 8 bước & Tổng số bước thực thi tối đa của các chuyên gia trong một lượt hội thoại\\
\texttt{MAX\_KNOWLEDGE\_RETRIEVALS} & 2 lần & Số lượt truy xuất tối đa vào kho tri thức vector PostgreSQL\\\bottomrule
\end{tabularx}
\caption{Các hằng số ràng buộc bất biến trong kiến trúc điều phối V2.}
\end{figure}

Quy trình diễn ra qua các giai đoạn:
\begin{enumerate}
    \item **Định tuyến ý định (Intent Routing):** Phân tích câu hỏi để xác định các chuyên gia cần kích hoạt. Nếu câu hỏi mơ hồ hoặc chứa các điều kiện mâu thuẫn (ví dụ: ``giá lớn hơn 200.000đ nhưng nhỏ hơn 100.000đ''), Supervisor lập tức trả về thông điệp yêu cầu làm rõ mà không lãng phí tài nguyên gọi các chuyên gia.
    \item **Lọc ứng viên:** Chuyên gia \texttt{product} thực thi tìm kiếm và chọn ra tối đa 5 ứng viên phù hợp nhất.
    \item **Thực thi chuyên gia song song:** Với tập ứng viên đã khóa, các chuyên gia \texttt{review} và \texttt{trust} được triệu hồi để thu thập chứng cứ.
    \item **Vòng lặp thích ứng theo dữ chứng (Adaptive Continuation Loop):** Nếu kết quả chặng 1 cho thấy cần làm rõ thêm thông tin ấn bản hoặc truy xuất nguồn gốc tác phẩm, Supervisor được phép thực hiện tối đa 1 lần bổ sung truy vấn.
    \item **Tổng hợp có kiểm chứng (Grounded Synthesis):** Ghép nối các mảnh dữ chứng thành câu trả lời hoàn chỉnh.
\end{enumerate}

\subsection{Cơ chế Đảm bảo Căn cứ và Trích dẫn Chuẩn hóa}
Điểm phân biệt căn bản giữa Evidence Atlas và các chatbot thông thường là **Sự sở hữu dữ chứng thuộc về máy chủ Python (Python-owned Provenance)**. Mô hình ngôn ngữ không được phép tự tạo ra mã dữ chứng hay bịa đặt liên kết nguồn \src{1}.

\subsubsection{Cơ chế Kiểm tra Xâu Ký tự Chính xác (Exact-Span Verification)}
Khi mô hình tổng hợp câu trả lời và trích dẫn một câu nhận xét từ người mua hoặc một thông tin từ kho tri thức, hệ thống thực thi thuật toán đối soát xâu văn bản chính xác:
\begin{equation}
\text{Match}(S_{\text{claim}}, D_{\text{source}}) = 
\begin{cases} 
1 & \text{nếu } S_{\text{claim}} \subseteq D_{\text{source}} \\
0 & \text{nếu } S_{\text{claim}} \not\subseteq D_{\text{source}}
\end{cases}
\end{equation}
Nếu phát hiện mô hình tự ý chỉnh sửa nội dung trích dẫn hoặc thêm thắt các chi tiết không có trong văn bản nguồn $D_{\text{source}}$, hệ thống kích hoạt cơ chế sửa chữa bản thảo (\texttt{Draft Repair}) hoặc từ chối công bố đoạn trích dẫn đó, thay thế bằng phát biểu trung tính có kiểm chứng.

\subsubsection{Định dạng Trích dẫn Thống nhất}
Mọi dữ chứng hợp lệ được hiển thị trong câu trả lời theo khuôn mẫu định dạng thống nhất `[C1]`, `[C2]`, tương ứng với hồ sơ dữ chứng (\textit{Evidence Dossier}) được trả về kèm theo trong phản hồi API. Người dùng có thể nhấp vào từng mã trích dẫn để xem chính xác văn bản gốc, nguồn trích lục và thời điểm ghi nhận.

\subsection{Quản lý Ngân sách và Giám sát Tài nguyên}
Để ngăn chặn nguy cơ cạn kiệt tài khoản API (\textit{denial-of-wallet}), hệ thống tích hợp mô-đun quản lý hạn mức ngân sách \texttt{ProviderBudgetScope} với độ phân giải siêu nhỏ đến cấp nano USD ($10^{-9}$ USD) \src{1}.

Mỗi lượt hội thoại được cấp một ngân sách tối đa định trước (mặc định tương đương $0{,}05$ USD). Trong quá trình xử lý, sau mỗi lần gọi mô hình ngôn ngữ hoặc tính toán vector embedding, hệ thống tự động trừ dần chi phí thực tế dựa trên số lượng prompt tokens và completion tokens đã sử dụng. 

Nếu phát hiện chi phí chạm ngưỡng ngân sách hoặc thời gian xử lý vượt quá thời hạn (\textit{timeout}), Supervisor lập tức hủy bỏ luồng thực thi và chuyển sang trạng thái ngắt an toàn (\texttt{BudgetCancelledError}), bảo vệ hệ thống khỏi các vòng lặp ngoài dự kiến.

\subsection{Bảo toàn Trạng thái Bền vững và Khóa Tranh chấp SQL}
Evidence Atlas lưu trữ toàn bộ lịch sử hội thoại, các lượt xử lý và các bản ghi hành động trong các bảng cơ sở dữ liệu PostgreSQL bền vững. 

Để giải quyết triệt để xung đột đồng thời (\textit{race condition}) khi máy khách gửi liên tiếp nhiều câu hỏi trong cùng một cuộc trò chuyện, hệ thống sử dụng cơ chế **Khóa dòng cơ sở dữ liệu (SQL Lease Lock)** thông qua câu lệnh:
\begin{verbatim}
SELECT * FROM conversations WHERE id = :conv_id FOR UPDATE NOWAIT;
\end{verbatim}
Cơ chế này đảm bảo rằng tại một thời điểm, chỉ duy nhất một tiến trình được phép xử lý lượt trò chuyện của cuộc hội thoại đó. Nếu một yêu cầu thứ hai đến khi lượt trước chưa hoàn thành, hệ thống sẽ trả về lỗi xung đột rõ ràng thay vì làm hỏng trạng thái hội thoại, loại bỏ sự phức tạp và rủi ro mất mát dữ liệu của các bộ điều phối phân tán trên bộ nhớ tạm.
